\documentclass[a4paper,12pt,twoside]{memoir}

% Castellano
\usepackage[spanish,es-tabla]{babel}
\selectlanguage{spanish}
\usepackage[utf8]{inputenc}
\usepackage[T1]{fontenc}
\usepackage{lmodern} % scalable font
\usepackage{microtype}
\usepackage{placeins}

\RequirePackage{booktabs}
\RequirePackage[table]{xcolor}
\RequirePackage{xtab}
\RequirePackage{multirow}

% Links
\PassOptionsToPackage{hyphens}{url}\usepackage[colorlinks]{hyperref}
\hypersetup{
	allcolors = {red}
}

% Ecuaciones
\usepackage{amsmath}

% Rutas de fichero / paquete
\newcommand{\ruta}[1]{{\sffamily #1}}

% Párrafos
\nonzeroparskip

% Huérfanas y viudas
\widowpenalty100000
\clubpenalty100000

% Evitar solapes en el header
\nouppercaseheads

% Imagenes
\usepackage{graphicx}
\newcommand{\imagen}[3]{
	\begin{figure}[!h]
		\centering
		\includegraphics[width=#3\textwidth]{#1}
		\caption{#2}\label{fig:#1}
	\end{figure}
	\FloatBarrier
}

\newcommand{\imagenflotante}[2]{
	\begin{figure}%[!h]
		\centering
		\includegraphics[width=0.9\textwidth]{#1}
		\caption{#2}\label{fig:#1}
	\end{figure}
}



% El comando \figura nos permite insertar figuras comodamente, y utilizando
% siempre el mismo formato. Los parametros son:
% 1 -> Porcentaje del ancho de página que ocupará la figura (de 0 a 1)
% 2 --> Fichero de la imagen
% 3 --> Texto a pie de imagen
% 4 --> Etiqueta (label) para referencias
% 5 --> Opciones que queramos pasarle al \includegraphics
% 6 --> Opciones de posicionamiento a pasarle a \begin{figure}
\newcommand{\figuraConPosicion}[6]{%
  \setlength{\anchoFloat}{#1\textwidth}%
  \addtolength{\anchoFloat}{-4\fboxsep}%
  \setlength{\anchoFigura}{\anchoFloat}%
  \begin{figure}[#6]
    \begin{center}%
      \Ovalbox{%
        \begin{minipage}{\anchoFloat}%
          \begin{center}%
            \includegraphics[width=\anchoFigura,#5]{#2}%
            \caption{#3}%
            \label{#4}%
          \end{center}%
        \end{minipage}
      }%
    \end{center}%
  \end{figure}%
}

%
% Comando para incluir imágenes en formato apaisado (sin marco).
\newcommand{\figuraApaisadaSinMarco}[5]{%
  \begin{figure}%
    \begin{center}%
    \includegraphics[angle=90,height=#1\textheight,#5]{#2}%
    \caption{#3}%
    \label{#4}%
    \end{center}%
  \end{figure}%
}
% Para las tablas
\newcommand{\otoprule}{\midrule [\heavyrulewidth]}
%
% Nuevo comando para tablas pequeñas (menos de una página).
\newcommand{\tablaSmall}[5]{%
 \begin{table}
  \begin{center}
   \rowcolors {2}{gray!35}{}
   \begin{tabular}{#2}
    \toprule
    #4
    \otoprule
    #5
    \bottomrule
   \end{tabular}
   \caption{#1}
   \label{tabla:#3}
  \end{center}
 \end{table}
}

%
%Para el float H de tablaSmallSinColores
\usepackage{float}

%
% Nuevo comando para tablas pequeñas (menos de una página).
\newcommand{\tablaSmallSinColores}[5]{%
 \begin{table}[H]
  \begin{center}
   \begin{tabular}{#2}
    \toprule
    #4
    \otoprule
    #5
    \bottomrule
   \end{tabular}
   \caption{#1}
   \label{tabla:#3}
  \end{center}
 \end{table}
}

\newcommand{\tablaApaisadaSmall}[5]{%
\begin{landscape}
  \begin{table}
   \begin{center}
    \rowcolors {2}{gray!35}{}
    \begin{tabular}{#2}
     \toprule
     #4
     \otoprule
     #5
     \bottomrule
    \end{tabular}
    \caption{#1}
    \label{tabla:#3}
   \end{center}
  \end{table}
\end{landscape}
}

%
% Nuevo comando para tablas grandes con cabecera y filas alternas coloreadas en gris.
\newcommand{\tabla}[6]{%
  \begin{center}
    \tablefirsthead{
      \toprule
      #5
      \otoprule
    }
    \tablehead{
      \multicolumn{#3}{l}{\small\sl continúa desde la página anterior}\\
      \toprule
      #5
      \otoprule
    }
    \tabletail{
      \hline
      \multicolumn{#3}{r}{\small\sl continúa en la página siguiente}\\
    }
    \tablelasttail{
      \hline
    }
    \bottomcaption{#1}
    \rowcolors {2}{gray!35}{}
    \begin{xtabular}{#2}
      #6
      \bottomrule
    \end{xtabular}
    \label{tabla:#4}
  \end{center}
}

%
% Nuevo comando para tablas grandes con cabecera.
\newcommand{\tablaSinColores}[6]{%
  \begin{center}
    \tablefirsthead{
      \toprule
      #5
      \otoprule
    }
    \tablehead{
      \multicolumn{#3}{l}{\small\sl continúa desde la página anterior}\\
      \toprule
      #5
      \otoprule
    }
    \tabletail{
      \hline
      \multicolumn{#3}{r}{\small\sl continúa en la página siguiente}\\
    }
    \tablelasttail{
      \hline
    }
    \bottomcaption{#1}
    \begin{xtabular}{#2}
      #6
      \bottomrule
    \end{xtabular}
    \label{tabla:#4}
  \end{center}
}

%
% Nuevo comando para tablas grandes sin cabecera.
\newcommand{\tablaSinCabecera}[5]{%
  \begin{center}
    \tablefirsthead{
      \toprule
    }
    \tablehead{
      \multicolumn{#3}{l}{\small\sl continúa desde la página anterior}\\
      \hline
    }
    \tabletail{
      \hline
      \multicolumn{#3}{r}{\small\sl continúa en la página siguiente}\\
    }
    \tablelasttail{
      \hline
    }
    \bottomcaption{#1}
  \begin{xtabular}{#2}
    #5
   \bottomrule
  \end{xtabular}
  \label{tabla:#4}
  \end{center}
}



\definecolor{cgoLight}{HTML}{EEEEEE}
\definecolor{cgoExtralight}{HTML}{FFFFFF}

%
% Nuevo comando para tablas grandes sin cabecera.
\newcommand{\tablaSinCabeceraConBandas}[5]{%
  \begin{center}
    \tablefirsthead{
      \toprule
    }
    \tablehead{
      \multicolumn{#3}{l}{\small\sl continúa desde la página anterior}\\
      \hline
    }
    \tabletail{
      \hline
      \multicolumn{#3}{r}{\small\sl continúa en la página siguiente}\\
    }
    \tablelasttail{
      \hline
    }
    \bottomcaption{#1}
    \rowcolors[]{1}{cgoExtralight}{cgoLight}

  \begin{xtabular}{#2}
    #5
   \bottomrule
  \end{xtabular}
  \label{tabla:#4}
  \end{center}
}




\graphicspath{ {./img/} }

% Capítulos
\chapterstyle{bianchi}
\newcommand{\capitulo}[2]{
	\setcounter{chapter}{#1}
	\setcounter{section}{0}
	\setcounter{figure}{0}
	\setcounter{table}{0}
	\chapter*{#2}
	\addcontentsline{toc}{chapter}{#2}
	\markboth{#2}{#2}
}

% Apéndices
\renewcommand{\appendixname}{Apéndice}
\renewcommand*\cftappendixname{\appendixname}

\newcommand{\apendice}[1]{
	%\renewcommand{\thechapter}{A}
	\chapter{#1}
}

\renewcommand*\cftappendixname{\appendixname\ }

% Formato de portada
\makeatletter
\usepackage{xcolor}
\newcommand{\tutor}[1]{\def\@tutor{#1}}
\newcommand{\course}[1]{\def\@course{#1}}
\definecolor{cpardoBox}{HTML}{E6E6FF}
\def\maketitle{
  \null
  \thispagestyle{empty}
  % Cabecera ----------------
\noindent\includegraphics[width=\textwidth]{cabecera}\vspace{1cm}%
  \vfill
  % Título proyecto y escudo informática ----------------
  \colorbox{cpardoBox}{%
    \begin{minipage}{.8\textwidth}
      \vspace{.5cm}\Large
      \begin{center}
      \textbf{TFG del Grado en Ingeniería Informática}\vspace{.6cm}\\
      \textbf{\LARGE\@title{}}
      \end{center}
      \vspace{.2cm}
    \end{minipage}

  }%
  \hfill\begin{minipage}{.20\textwidth}
    \includegraphics[width=\textwidth]{escudoInfor}
  \end{minipage}
  \vfill
  % Datos de alumno, curso y tutores ------------------
  \begin{center}%
  {%
    \noindent\LARGE
    Presentado por \@author{}\\ 
    en Universidad de Burgos --- \@date{}\\
    Tutores: \@tutor{}\\
  }%
  \end{center}%
  \null
  \cleardoublepage
  }
\makeatother


% Datos de portada
\title{Comparación de modelos de predicción de datos de temperatura y humedad procedentes de sensores IoT desplegados en un viñedo \\Documentación Técnica}
\author{Junhao Zhou}
\tutor{Alejandro Merino Gómez, Rubén Ruiz González}
\date{\today}

\begin{document}

\maketitle



\cleardoublepage



%%%%%%%%%%%%%%%%%%%%%%%%%%%%%%%%%%%%%%%%%%%%%%%%%%%%%%%%%%%%%%%%%%%%%%%%%%%%%%%%%%%%%%%%



\frontmatter


\clearpage

% Indices
\tableofcontents

\clearpage

\listoffigures

\clearpage

\listoftables

\clearpage

\mainmatter

\appendix

\apendice{Plan de Proyecto Software}

\section{Introducción}

En este apéndice se explicarán la planificación realizada para el desarrollo del proyecto y la viabilidad de esta.

\section{Planificación temporal}

La planificación del proyecto y su desarrollo se han realizado siguiendo una metodología ágil, en este caso se ha empleado Scrum. Scrum emplea ciclos de desarrollo o iteraciones (\textit{sprints}), que definen cada una de las etapas de trabajo con un objetivo concreto dentro del desarrollo. Se ha intentado que la duración de los \textit{sprints} fueran consistentes, y sean de una duración de una o dos semanas.

\subsection{01 - Preprocesamiento de datos [03/10/2022 - 31/10/2022]}

Este primer \textit{sprint} se han realizado las siguientes tareas:

\begin{itemize}
    \item Estudio del proceso CRISP-DM: se ha estudiado el proceso y sus fases para aplicarlo al proyecto.
    \item Estudio de la biblioteca de Python Pandas: dado que es la primera biblioteca que se va a emplear para poder importar y emplear el conjunto de datos en Python, se ha realizado un estudio de la biblioteca y de las funciones que tiene para el manejo de datos.
    \item Validación de los datos: se ha intentado realizar una validación de los datos para descartar aquellos datos que son erróneos.
\end{itemize}

\subsection{02 - Validación de datos [31/10/2022 - 14/11/2022]}

En este \textit{sprint} se han realizado las tareas mencionadas a continuación:

\begin{itemize}
    \item Creación de columnas de validación: se han creado una columna que contiene la información sobre si ese dato es válido en todos sus campos empleando una cadena numérica.
    \item Descarga de datos de una API: se han obtenido datos meteorológicos históricos para validar los datos comparándolos del conjunto de datos con los obtenidos por la API.
    \item Validación de los datos: se ha continuado con la validación de los datos, para ello los comparamos con los de la API.
    \item Observación de relaciones con los datos válidos: una vez validamos los datos, queremos realizar un análisis de correlación lineal para ver las relaciones entre los atributos, como, por ejemplo, si la humedad sube cuando hay lluvia o la humedad disminuye cuando sube la temperatura.
\end{itemize}

\subsection{03 - Investigación para validar datos [14/11/2022 - 05/12/2022]}

Durante este \textit{sprint} se han hecho las siguientes tareas:

\begin{itemize}
    \item Creación del repositorio en Github: se ha creado el repositorio en Github para poder gestionar los cambios en el código.
    \item Investigación para validar datos: se ha investigado acerca de los márgenes adecuados a aplicar para descartar los datos, ya que queremos descartar el menor número de datos posible y que todos ellos sean válidos.
\end{itemize}

\subsection{04 - Continuación en la validación de datos [05/12/2022 - 30/01/2023]}

En este \textit{sprint} se ha continuado con la validación de datos, ya que al estar trabajando con un conjunto de datos sin tratar, tenemos que realizar todo el proceso de preprocesamiento, incluyendo la investigación y las pruebas en la validación de datos.

\subsection{05 - Cambio en la validación de datos [30/01/2023 - 20/02/2023]}

Durante este \textit{sprint} se ha seguido con la validación de los datos. Se ha eliminado la columna que contenía la información sobre la validez del dato y se han añadido tantas columnas de validez como atributos queremos validar en cada dato. Se ha empleado una codificación ternaria en dichas columnas para indicar si el atributo es válido, no válido o no se ha validado.

\subsection{06 - Cambio de API [20/02/2023 - 27/02/2023]}

En ese \textit{sprint} se ha cambiado la API empleada y se ha usado Weatherbit en vez de OpenWeatherMap y se han hecho las llamadas correspondientes para volver a obtener los datos meteorológicos históricos.

\subsection{07 - Cambio en la manera de validar [27/02/2023 - 06/03/2023]}

En este \textit{sprint} se ha modificado la manera en la que se validan los datos, puesto que se descartaban demasiados datos con el anterior método de validación. Ya no usamos los datos de la API para compararlos con los de los sensores, en cambio, empleamos una media móvil de los datos de todos los sensores para validar cada uno de los datos de cada sensor.

\subsection{08 - Análisis exploratorio de datos [06/03/2023 - 13/03/2023]}

En este \textit{sprint} se ha realizado un análisis exploratorio del conjunto de datos y un estudio de correlaciones lineales de las variables para observar las posibles relaciones que pueda haber entre ellas.

\subsection{09 - Adición de columnas [13/03/2023 - 20/03/2023]}

En este \textit{sprint} se han añadido columnas de variación de la temperatura y la humedad exterior al conjunto de datos.

\subsection{10 - Investigación y desarrollo de modelos para predicciones de un solo paso [20/03/2023 - 10/04/2023]}

En este \textit{sprint} se han realizado las siguientes tareas:

\begin{itemize}
    \item Investigación de modelos de predicción de series temporales: se han investigado y se ha realizado una comparativa inicial de los modelos de predicción de un solo paso para escoger qué modelos desarrollar.
    \item Desarrollo de modelos de predicción de un solo paso: se ha empezado el desarrollo de los modelos de predicción de un solo paso escogidos en base a la investigación anterior.
    \item Adición de columna de meses: se ha añadido una columna para indicar el mes del año en el que se ha recogido el dato para proporcionar más información a los modelos a entrenar.
\end{itemize}

\subsection{11 - Desarrollo de modelos de predicción de un solo paso [10/04/2023 - 24/04/2023]}

En este \textit{sprint} se han desarrollado y modificado los modelos lineal y LSTM empleados para las predicciones de un solo paso para conseguir la mejor precisión posible y se ha realizado una comparativa del error devuelto por cada modelo.

\subsection{12 - Investigación y desarrollo de modelos para predicciones de múltiples pasos [24/04/2023 - 16/05/2023]}

En este \textit{sprint} se han realizado las siguientes tareas:

\begin{itemize}
    \item Investigación de modelos de predicción de múltiples pasos: se ha investigado acerca de las redes neuronales que se pueden aplicar para predicciones de múltiples pasos.
    \item Inicio del desarrollo de los modelos: se ha empezado el desarrollo de las redes neuronales que se van a emplear para este tipo de predicciones.
    \item Cambio de normalización: se ha cambiado la técnica de normalización empleada, pasando de la puntuación estándar a la normalización Min-Max.
    \item Adición del atributo hora: se ha añadido al conjunto de datos empleado el atributo hora para que los modelos sean capaces de realizar predicciones con mayor precisión.
\end{itemize}

\subsection{13 - Optimización de los modelos y modificación del conjunto de datos [16/05/2023 - 29/05/2023]}

En este \textit{sprint} se han hecho las siguientes tareas:

\begin{itemize}
    \item Modificación del conjunto de datos: se quiere escoger un periodo en el que no haya saltos en los datos, para que todas las subsecciones creadas a partir del conjunto disponga de un subconjunto de datos donde todos los pasos de tiempo son continuos. Esto se realiza con la finalidad de optimizar los resultados de los modelos.
    \item Optimización de los modelos: se han realizado varias pruebas con distintos parámetros de configuración de los modelos para encontrar los más óptimos.
\end{itemize}

\subsection{14 - Adición de nuevos modelos [29/05/2023 - 06/06/2023]}

En este \textit{sprint} se han añadido dos nuevos modelos para predicciones de múltiples pasos, que son el modelo convolucional y el modelo \textit{Dense}.

\subsection{15 - Redacción de la memoria y los anexos [06/06/2023 - 22/06/2023]}

En este penúltimo \textit{sprint} se ha redactado la memoria y los anexos, siendo los últimos procesos a realizar del proyecto.

\subsection{16 - Retoques finales [22/06/2023 - 30/06/2023]}

En el último \textit{sprint} del proyecto se han realizado las siguientes tareas:

\begin{itemize}
    \item Finalización de la memoria: se ha terminado de redactar la memoria del proyecto.
    \item Modificación del código: se ha modificado el fichero que mostraba las gráficas para añadir imágenes a la memoria.
    \item Finalización de los anexos: se han terminado de redactar los distintos anexos.
\end{itemize}

\section{Estudio de viabilidad}

A continuación se realizará un breve estudio de la viabilidad económica y legal del proyecto.

\subsection{Viabilidad económica}

Este proyecto está solamente enfocado al desarrollo y uso de modelos, y no se ha planeado obtener ningún beneficio económico durante su desarrollo, por lo que no se podría considerar como un producto para publicarlo en el mercado. Además, tampoco se dispone de una aplicación que podamos publicar que permita la ejecución de los modelos.

\par

Sin embargo, aunque no tiene como objetivo obtener ganancias económicas, es importante considerar los coses asociados al proyecto.

\par

Algunos aspectos a tener en cuenta para realizar los cálculos son los siguientes:

\begin{itemize}
    \item Suponemos que el proyecto se realiza en una empresa ubicada en España.
    \item El proyecto tiene una duración aproximada de 9 meses, desde principios de octubre hasta principios de julio. Lo que equivale a 36 semanas aproximadamente, considerando una duración de 4 semanas por mes.
    \item El equipo encargado del proyecto está compuesto por un desarrollador (el alumno), el \textit{Product Owner} y el \textit{Scrum Master}, que son los tutores del alumno en el proyecto.
\end{itemize}

\subsubsection{Costes}

Para empezar, empezaremos calculando los costes de la empresa.

\begin{enumerate}
    \item \textbf{Empleados}
    \par
    Calcularemos el salario de cada uno de los empleados teniendo en cuenta el salario promedio en España.
    \begin{itemize}
        \item \textbf{Desarrollador}: dado que el desarrollador es considera un programador \textit{junior}, su salario bruto en España sería de 19.000 € anuales~\cite{salario:junior}. Teniendo en cuenta la duración del proyecto y suponiendo que las pagas mensuales estan prorrateadas, podemos sacar el coste total con la siguiente fórmula:
        
        \[\textup{Salario bruto mensual} = \frac{\textup{Salario bruto anual}}{\textup{12 pagas}}\]

        \par

        Sustituyendo los valores, obtenemos el coste mensual, que es de:

        \[\textup{Salario bruto mensual} = \frac{\textup{19\,700 €}}{\textup{12 pagas}} \approx \textup{1641,67 €}\]
        \item \textbf{Product Owner}: el salario medio es de 45.000 € brutos anuales~\cite{salario:owner}. El salario mensual de este empleado es de:

        \[\textup{Salario bruto mensual} = \frac{\textup{45\,000 €}}{\textup{12 pagas}} \approx \textup{3750 €}\]

        \item \textbf{Scrum Master}: el salario medio de este empleado es el mismo que el del \textit{Product Owner}~\cite{salario:master}.
        
    \end{itemize}

    Por último, también tenemos que tener en cuenta las horas aplicadas al proyecto por parte del \textit{Product Owner} y del \textit{Scrum Master} para calcular el coste real del proyecto. Vamos a suponer que dedican un 5\% de su jornada laboral al proyecto, por lo que, los costes de los salarios que se aplican al proyecto son de:

    \[\textup{Coste por empleado} = \textup{3.750 €} \times 0,05 \approx \textup{187,5 €}\]

    \par

    Teniendo en cuenta los salarios brutos, podemos calcular el coste total de los empleados sin tener en cuenta los costes adicionales a pagar por la empresa por la Seguridad Social.

    \par

    El coste total de los salarios brutos es de 18.150,03 € aproximadamente.
    
    \item\textit{\textbf{Hardware}}
    \par

    Como costes de \textit{hardware} incluiremos el ordenador del alumno, que ha sido lo que se ha empleado para el desarrollo del proyecto. El equipo dispone de un procesador Intel(R) Core(TM) i7-7700K CPU @ 4.20GHz y 16 GB de memoria RAM. Este equipo se compró hace 6 años a un precio de aproximadamente 1.300 €.

    Vamos a suponer la amortización del equipo, suponiendo una vida útil del equipo de 10 años. La amortización anual del equipo es de 130 €, por lo que para este proyecto necesitamos sacar el 75\% de ese coste.

    Los costes de \textit{hardware} de este proyecto han sido de 97,5 €.

    \item \textbf{\textit{Software}}
    \par

    Como costes de \textit{software} podemos incluir las licencias de las herramientas empleadas que tengan un coste. Ya que en el caso del proyecto se han empleado licencias temporales gratuitas, vamos a suponer los costes una vez finalicen los plazos de estas.

    \begin{itemize}
        \item \textbf{Weatherbit}: para obtener los datos de la API necesitaríamos una licencia <<Pro>> como mínimo, cuyo coste es de 180 € al mes. Ya que solo lo usaríamos el primer mes para obtener los datos históricos y ya, solo contabilizamos un mes de costes.
        \item \textbf{Github}: tendríamos que pagar la versión \textit{Enterprise} del Github, que tiene un coste de 19,64 € anuales aproximadamente.
        \item \textbf{MATLAB}: por último, como también se emplea MATLAB en una fase del proyecto, debemos tener en cuenta el coste de su licencia. En el caso supuesto, necesitaríamos la licencia \textit{Standard}. que tiene un coste anual de 860 €.
    \end{itemize}

    A partir de los gastos mencionados, podemos obtener el coste del proyecto, que es de:

    \[\textup{Gastos de software} = \textup{180 €} + (\textup{19,64 € + 860 €}) \times \textup{0,75} \approx \textup{837,48 €}\]
\end{enumerate}

En la tabla~\ref{tabla:costes} podemos observar los gastos totales del proyecto.

\tablaSmall{Costes totales del proyecto}{ l l l }{costes}{Costes & & Coste Total\\}
{Empleados & & 18.150,03 € €\\
Hardware & & 97,5 €\\
Software & & 837,48 €\\
\hline\hline
\textbf{Total} & & 19.085,01 €\\}

\subsection{Viabilidad legal}

Respecto a la viabilidad legal, podemos enfocarlo en dos puntos: las herramientas y los datos empleados.

\par

Respecto a las herramientas empleadas en este proyecto, es importante tener en cuenta las licencias de cada una de ellas. A continuación, se van a enumerar todas las licencias de las herramientas empleadas:

\begin{itemize}
    \item Licencia MIT: la licencia MIT permite el uso, copia, modificación y distribución del \textit{software}, tanto en proyectos comerciales como no comerciales, con la condición de mantener el aviso de derechos de autor y la exención de responsabilidad.
    \item Licencia BSD: esta licencia permite el uso, modificación y distribución del \textit{software}.
    \item Licencia Open Source: esta licencia permite su uso y distribución para uso personal y comercial.
    \item Licencia GNU: esta licencia permite el uso, estudio, modificación y distribución del \textit{software}. Además obliga a que todas las modificaciones y actualizaciones sigan siendo de código abierto.
    \item Licencia Apache: la licencia Apache permite el uso, modificación y distribución del \textit{software} siempre que se mantenga el aviso de derechos de autor, la exención de responsabilidad y la necesidad de proporcionar una copia de la licencia en cualquier distribución de la herramienta.
    \item Licencia Creative Commons: esta licencia permite el uso del \textit{software} a los que la posean sin solicitar el permiso del autor.
    \item Licencia Open Database: esta licencia permite a los usuarios el uso, modificación y distribución del \textit{software}.
    \item Licencia de uso comercial: esta licencia permite a cualquier empleado de una empresa que la tenga el uso del \textit{software}. Esta licencia es de pago. 
\end{itemize}

\par

En la tabla~\ref{tabla:licencias} podemos observar las herramientas y sus licencias.

\begin{center}
\begin{tabular}{ |l| l| } 
 \hline
 \textbf{Herramienta} & \textbf{Licencia} \\ 
 \hline
 Visual Studio Code & MIT License  \\ 
 \hline
 Python & Open source License  \\ 
 \hline
 MATLAB & Commercial Use License\\
 \hline
 matplotlib & BSD License  \\ 
 \hline
 seaborn & BSD License\\
 \hline
 NumPy & BSD License\\
 \hline
 Pandas & BSD License  \\ 
 \hline
 TensorFlow & Open source License  \\ 
 \hline
 OpenWeatherMap & OpenWeather end-user License~\cite{licencia:OpenWeather}\\
 \hline
 Weatherbit & Commercial Use License\\
 \hline
\end{tabular}
\label{tabla:licencias}
\end{center}

\par

En cuanto a los datos empleados en el proyecto, puesto que el conjunto de datos empleado es de propiedad privada, solo es válido para la investigación y el desarrollo de proyectos que permita el dueño de estos. Por lo tanto, los modelos entrenados en este proyecto solo se pueden emplear para el viñedo de donde proceden los datos.
\apendice{Especificación de Requisitos}

\section{Introducción}

Al estar empleando \textit{Scrum} como metodología ágil, se van a describir a continuación las historias de usuario (requisitos) del proyecto.

\section{Objetivos generales}

Este proyecto cuenta con los siguientes objetivos técnicos:

\begin{itemize}
    \item Realizar predicciones de variables de interés agronómico en un viñedo.
    \item Usar regresión mediante redes neuronales para la predicción en series temporales.
    \item Realizar minería de datos sobre el conjunto de datos propuesto para tratarlos previamente al desarrollo de redes neuronales.
    \item Comprender el funcionamiento de los distintos modelos de predicción a desarrollar.
    \item Desarrollar múltiples modelos de predicción empleando distintos tipos de redes neuronales.
    \item Realizar un análisis comparativo de los modelos desarrollados con los modelos de otros proyectos relacionados.
\end{itemize}

\section{Catálogo de requisitos}

A partir de los objetivos descritos en el apartado anterior, podemos desarrollar las siguientes historias de usuario o \textit{User Stories} (US):

\begin{itemize}
    \item \textbf{HU-01} El usuario debe ser capaz de cargar los datos.
    \item \textbf{HU-02} El usuario debe ser capaz de preparar los datos para el entrenamiento de los modelos.
    \item \textbf{HU-03} El usuario debe ser capaz de modificar los parámetros de entrenamiento de los modelos.
    \item \textbf{HU-04} El usuario debe poder modificar los parámetros de configuración del modelo.
    \item \textbf{HU-05} El usuario debe ser capaz de entrenar un modelo.
    \item \textbf{HU-06} El usuario debe ser capaz de validar un modelo.
    \item \textbf{HU-07} El usuario debe ser capaz de visualizar los resultados.
\end{itemize}

\section{Especificación de requisitos}

\subsection{Diagrama de casos de uso}

\imagen{diagrama-cu}{Diagrama de casos de uso}{1}

\subsection{Especificaciones de los casos de uso}

% Caso de Uso 1 -> Carga de datos.
\begin{table}[p]
	\centering
	\begin{tabularx}{\linewidth}{ p{0.21\columnwidth} p{0.71\columnwidth} }
		\toprule
		\textbf{CU-1}    & \textbf{Carga de datos}\\
		\toprule
		\textbf{Versión}              & 1.0    \\
		\textbf{Requisitos asociados} & HU-01 \\
		\textbf{Descripción}          & El usuario puede cargar los datos que se van a emplear para el entrenamiento de los modelos \\
		\textbf{Precondición}         & El usuario debe disponer de un conjunto de datos \\
		\textbf{Acciones}             &
		\begin{enumerate}
			\def\labelenumi{\arabic{enumi}.}
			\tightlist
			\item El usuario carga los datos del fichero en el código.
		\end{enumerate}\\
		\textbf{Postcondición}        & Se debe haber almacenado el conjunto de datos a emplear en una variable. \\
		\textbf{Excepciones}          & El fichero tiene una extensión no aceptada. \\
		\textbf{Importancia}          & Alta \\
		\bottomrule
	\end{tabularx}
	\caption{CU-1 Carga de datos.}
\end{table}

% Caso de Uso 2 -> Preparación de datos.
\begin{table}[p]
	\centering
	\begin{tabularx}{\linewidth}{ p{0.21\columnwidth} p{0.71\columnwidth} }
		\toprule
		\textbf{CU-2}    & \textbf{Preparación de datos}\\
		\toprule
		\textbf{Versión}              & 1.0    \\
		\textbf{Requisitos asociados} & HU-02 \\
		\textbf{Descripción}          & El usuario puede modificar el conjunto de datos cargado para emplear ciertas variables en vez de todas y descartar los datos que no se quieran emplear. También debe ser capaz de normalizar, dividir el conjunto de datos en \textit{sets} y generar secciones de datos. \\
		\textbf{Precondición}         & El usuario debe disponer de un conjunto de datos. El usuario debe poder cargar datos. \\
		\textbf{Acciones}             &
		\begin{enumerate}
			\def\labelenumi{\arabic{enumi}.}
			\tightlist
			\item El usuario carga los datos.
            \item Si el usuario quiere descartar datos, descarta todos aquellos que no quiera usar para el entrenamiento.
            \item Si el usuario quiere descartar variables del conjunto de datos, descarta aquellos que no quiera emplear para el entrenamiento.
            \item El usuario normaliza los datos con la técnica deseada.
            \item El usuario divide el conjunto de datos en conjunto de entrenamiento, de validación y de test.
            \item El usuario genera secciones de datos para las predicciones de series temporales.
		\end{enumerate}\\
		\textbf{Postcondición}        & Se debe tener un conjunto de datos preparado para el entrenamiento de los modelos. \\
		\textbf{Excepciones}          & El conjunto de datos está vacío. La suma de las proporciones de los \textit{sets} creados supera el 100\%. \\
		\textbf{Importancia}          & Alta \\
		\bottomrule
	\end{tabularx}
	\caption{CU-2 Preparación de datos.}
\end{table}

% Caso de Uso 3 -> Entrenamiento y validación de modelos.
\begin{table}[p]
	\centering
	\begin{tabularx}{\linewidth}{ p{0.21\columnwidth} p{0.71\columnwidth} }
		\toprule
		\textbf{CU-3}    & \textbf{Entrenamiento y validación de modelos}\\
		\toprule
		\textbf{Versión}              & 1.0    \\
		\textbf{Requisitos asociados} & HU-03, HU-04, HU-05, HU-06 \\
		\textbf{Descripción}          & El usuario puede entrenar un modelo y validarlo a partir de un conjunto de datos preparado.\\
		\textbf{Precondición}         & El usuario debe disponer de un conjunto de datos preparado. El usuario debe disponer de un modelo para entrenar.\\
		\textbf{Acciones}             &
		\begin{enumerate}
			\def\labelenumi{\arabic{enumi}.}
			\tightlist
			\item El usuario modifica los parámetros del entrenamiento del modelo.
            \item El usuario modifica los parámetros de la configuración del modelo.
            \item El usuario entrena el modelo.
            \item El usuario valida el modelo.
		\end{enumerate}\\
		\textbf{Postcondición}        & Se debe haber obtenido un modelo entrenado y validado. Se debe haber obtenido el error de ese modelo. \\
		\textbf{Excepciones}          & El conjunto de datos está vacío. El modelo no está bien configurado.  \\
		\textbf{Importancia}          & Alta \\
		\bottomrule
	\end{tabularx}
	\caption{CU-3 Entrenamiento y validación de modelos.}
\end{table}

% Caso de Uso 4 -> Visualizar resultados.
\begin{table}[p]
	\centering
	\begin{tabularx}{\linewidth}{ p{0.21\columnwidth} p{0.71\columnwidth} }
		\toprule
		\textbf{CU-4}    & \textbf{Visualizar los resultados}\\
		\toprule
		\textbf{Versión}              & 1.0    \\
		\textbf{Requisitos asociados} & HU-07 \\
		\textbf{Descripción}          & El usuario puede visualizar el error de cada modelo en una gráfica.\\
		\textbf{Precondición}         & El usuario debe disponer de mínimo un modelo entrenado y validado.\\
		\textbf{Acciones}             &
		\begin{enumerate}
			\def\labelenumi{\arabic{enumi}.}
			\tightlist
			\item El usuario muestra el error de cada modelo en una gráfica.
		\end{enumerate}\\
		\textbf{Postcondición}        & Se debe haber obtenido una gráfica donde se represente el error de cada modelo. \\
		\textbf{Excepciones}          & No se ha entrenado ningún modelo.  \\
		\textbf{Importancia}          & Alta \\
		\bottomrule
	\end{tabularx}
	\caption{CU-4 Visualizar resultados.}
\end{table}
\apendice{Especificación de diseño}

\section{Introducción}

En este anexo se va a hablar del diseño de distintos apartados del proyecto.

\section{Diseño de datos}

A la hora de entrenar los distintos modelos de predicción de series temporales, es importante conocer la estructura del conjunto de datos a emplear. Dado que el conjunto de datos no es público, no se puede encontrar en ningún lugar a menos que se obtenga con permiso del propietario.

\par

Dado que se ha modificado el conjunto de datos a emplear para el entrenamiento de los modelos, se va a comentar el diseño del conjunto de datos final y sus características. El conjunto de datos empleado para el entrenamiento de los modelos tiene las siguientes características: 

\begin{itemize}
    \item 2987 datos u observaciones
    \item 6 columnas o atributos
    \item Los atributos son el mes, la hora, la temperatura y la humedad a baja y alta profundidad
\end{itemize}

A este conjunto de datos se le ha realizado un pequeño análisis exploratorio de datos, empleando una función proporcionada por la biblioteca de Python \textit{Pandas}. En la figura~\ref{fig:analisis_datos} se pueden observar algunos datos estadísticos del conjunto de datos empleado.

\imagen{analisis_datos}{Datos estadísticos del conjunto de datos empleado}{1}

Respecto al conjunto de datos de donde proviene el conjunto de datos empleado, es del archivo \texttt{SoloDatosValidos\_sensor3.csv}. Ese conjunto de datos tiene las siguientes características:

\begin{itemize}
    \item 68893 datos u observaciones
    \item 16 columnas o atributos
    \item Todos los atributos de este conjunto de datos son válidos y se pueden emplear todos los atributos.
\end{itemize}

Si se quisiera realizar un preprocesamiento de datos distinto, se recomienda realizarlo a partir del fichero \texttt{sensor3.csv}, que contiene todos los datos del sensor sin validar, o a partir del fichero \texttt{Datos\_Validados\_sensor3.csv}, que contiene todos los datos del sensor validados y con columnas adicionales indicando la validez de los datos.

\section{Diseño procedimental}

Para el desarrollo del proyecto se han realizado múltiples experimentos empleando distintos parámetros. Los distintos procedimientos realizados en este proyecto se pueden representar en diagramas de flujo donde se puede observar el flujo general seguido para estos.

En el diagrama de flujo de la figura~\ref{fig:procedimiento} se puede observar el procedimiento seguido para la preparación del conjunto de datos.

\imagen{procedimiento}{Diagrama de flujo de la preparación de los datos}{1}

El proceso seguido para la preparación del conjunto de datos es simple: se cargan los datos, se reduce el conjunto si se desea para trabajar con una menor cantidad de datos, pero asegurarnos de que no hay ningún salto en los pasos de tiempo, eliminar atributos del conjunto de datos si consideramos que no son útiles para las predicciones, normalizamos los datos, dividimos el conjunto de datos en \textit{sets} de entrenamiento, validación y test y generamos secciones de datos para realizar predicciones de series temporales.

Por otro lado, también es importante desarrollar el proceso seguido para el entrenamiento de los modelos, dicho proceso se puede ver en la figura~\ref{fig:proceso-entrenamiento}. Primero se han cargado los datos, se realiza el procedimiento para la preparación de los datos, se configuran los parámetros de entrenamiento para todos los modelos y se configuran los modelos. Si son modelos de un solo paso se entrena y se obtiene el error, en cambio, si son modelos de múltiples pasos se entrenan y se obtiene el error durante un número de iteraciones establecido en el código. Por último, calculamos la media de los errores en el caso de los modelos de múltiples pasos y mostramos los resultados para realizar la comparación.

\imagen{proceso-entrenamiento}{Diagrama de flujo del proceso seguido para el entrenamiento de los modelos}{1}

\section{Diseño arquitectónico}

La arquitectura de este proyecto es bastante simple, puesto que se centra en la investigación y la comparación de modelos de predicción. Debido a ello, se ha empleado un modelo arquitectónico llamado patrón pizarra~\cite{pizarra}.

\par

Este modelo arquitectónico consta de múltiples elementos funcionales, llamados agentes, y un instrumento de control, llamado pizarra. Los agentes están especializados en resolver una tarea concreta, en este caso, un entrenamiento de un modelo de predicción. Además, todos ellos cooperan para alcanzar un objetivo común, que en este proyecto sería la comparación de distintos modelos para predicciones de series temporales. La pizarra es el elemento que nos permite la visualización de todos los agentes y realizar operaciones con todos ellos. En este proyecto, el elemento pizarra sería el fichero Jupyter Notebook en el que se han entrenado los modelos y hemos visualizado los resultados.

\par

Podemos observar en la figura~\ref{fig:pizarra} el diseño arquitectónico del proyecto. Se puede ver como los agentes son cada uno de los modelos que se han entrenado y el conjunto de datos empleado y la pizarra es el fichero donde se ha realizado todo el proceso.

\imagen{pizarra}{Estructura arquitectónica del proyecto}{1}
\apendice{Documentación técnica de programación} \label{Documentacion_programacion}

\section{Introducción}

En esta sección se comentará la estructura de los directorios del proyecto y su proceso de ejecución para poder replicarlo.

\section{Estructura de directorios}

El proyecto tiene los siguientes directorios y ficheros:

\begin{itemize}
    \item /2022-10-01/: en este directorio se encuentran los datos recogidos por cada sensor sin tratar.
    \item /Datos API/: en este directorio se encuentran los datos meteorológicos históricos obtenidos a través de las APIs empleadas. Se encuentran los datos históricos de ambas APIs utilizadas.
    \item /DatosValidados/: en este directorio se encuentran los datos de cada sensor validados. Se encuentran los ficheros que contienen todos los datos con columnas de validación y los ficheros que solo contienen datos con todos los atributos válidos.
    \item /Graficas.ipynb: este fichero contiene código que devuelve gráficas de validez de los datos.
    \item /Predicciones.ipynb: este fichero contiene el código necesario para la preparación de los datos, la normalización de estos, la generación de secciones y el entrenamiento y validación de los modelos. También está la visualización de los resultados de cada modelo.
    \item /Validacion.ipynb: en este fichero se encuentra el código empleado para la obtención de datos a través de APIs y para la limpieza de los datos de los sensores. También se encuentra el código empleado para la adición de atributos al conjunto de datos.
\end{itemize}

\section{Compilación, instalación y ejecución del proyecto}

En esta sección se explicarán los pasos necesarios para replicar el entorno en el que se ha ejecutado el código del proyecto.

\subsection{Git}

Si queremos clonar el repositorio necesitamos tener Git instalado en el equipo. Para ello, podemos acceder a su página web~\cite{git}.

Una vez tenemos Git en el equipo, basta con clonar el repositorio. Para ello nos dirigimos al directorio donde queremos que se clone, abrimos una consola de comandos y escribimos el siguiente comando:

\texttt{git clone https://github.com/jzx1003/GII\_O\_MA\_22.07.git}

\subsection{Entorno de ejecución}

Para poder ejecutar el código, necesitamos un entorno donde ejecutar el código, en este proyecto se ha empleado Visual Studio Code~\cite{vscode} para la ejecución, aunque se puede realizar empleando solamente Anaconda~\cite{anaconda}. La versión de Python empleada es Python 3.10.9, por lo que si queremos emplear Anaconda para la ejecución, deberemos instalar la versión correspondiente que emplee esa versión de Python. No es necesario instalar Python aparte, puesto que el instalador de Anaconda también lo instala.

Tras la instalación de Anaconda, podemos emplear la herramienta Jupyter Notebook que trae Anaconda para abrir el fichero y ejecutarlo. 

Para su ejecución con Visual Studio Code deberemos instalar las extensiones de Jupyter ofrecidas por Microsoft en la pestaña de extensiones de la propia aplicación. Se pueden ver las extensiones en la figura~\ref{fig:vscode_jupyter}.

\imagen{vscode_jupyter}{Extensiones de Jupyter en Visual Studio Code}{1}

\section{Manual del programador}

En este último apartado se explicará cómo un programador puede ejecutar y modificar los parámetros del código. Para poder ejecutar el código necesitamos cumplir con lo mencionado en el apartado anterior.

\par

El fichero que contiene el código que entrena los distintos modelos es el fichero \texttt{Predicciones.ipynb}. El proceso que se sigue en ese fichero es el siguiente:

\begin{itemize}
    \item Carga de datos
    \item Preparación de los datos
    \item Normalización de los datos
    \item División del conjunto de datos
    \item Generación de secciones
    \item Configuración de los parámetros de entrenamiento
    \item Entrenamiento y validación de los modelos
    \item Visualización de los resultados
\end{itemize}

En la carga de datos, podemos modificar qué fichero queremos que cargue cambiando la ruta para que lea el fichero deseado y lo almacene en un DataFrame de Pandas. En la figura~\ref{fig:carga_datos} se puede observar la carga de datos.

\imagen{carga_datos}{Código de la carga de datos}{1}

En la preparación de datos podemos escoger el rango de los datos que queremos usar para el entrenamiento y los atributos que queremos descartar. En este caso, descartamos muchos de los atributos, pues se ha querido entrenar los modelos con una menor cantidad de datos y reducir la complejidad. En la figura~\ref{fig:preparacion_datos} se realiza la preparación del conjunto de datos.

\imagen{preparacion_datos}{Código de la preparación de los datos}{1}

A continuación, se realiza la normalización de los datos, en este caso se ha empleado la normalización Min-Max, se puede modificar el código para cambiar el método de normalización. En la figura~\ref{fig:normalizacion} se puede observar la celda empleada para la normalización.

\imagen{normalizacion}{Código de la normalización de los atributos}{1}

Para dividir el conjunto de datos en los \textit{sets} de entrenamiento, validación y test se ha empleado el código de la celda mostrada en la figura~\ref{fig:division_datos}. Se puede modificar el porcentaje de datos asignado a cada \textit{set} cambiando los valores decimales.

\imagen{division_datos}{Código de la división de datos}{1}

A continuación, podemos modificar las secciones generadas para las predicciones empleando la función WindowGenerator, creada para ello. Los parámetros a introducir en la función son los siguientes:

\begin{itemize}
    \item \texttt{input\_width}: Número de pasos de tiempo que introducimos como entrada
    \item \texttt{label\_width}: Número de pasos de tiempo a predecir
    \item \texttt{shift}: Desplazamiento de los pasos de tiempo a predecir respecto a los pasos de tiempo de entrada
    \item \texttt{label\_columns}: Nombre de los atributos a predecir y a mostrar
\end{itemize}

En la figura~\ref{fig:window_codigo} podemos observar las secciones creadas para las predicciones de un solo paso.

\imagen{window_codigo}{Código de las secciones para predicciones de un solo paso}{0.7}

También podemos modificar los parámetros de la función empleada para el entrenamiento y la validación de los modelos. La función creada se llama \texttt{compile\_and\_fit(model, window, patience)}. Introducimos como primer argumento el modelo a entrenar y validar, como segundo argumento la sección que queramos emplear para las predicciones y como último argumento la paciencia que queramos que tenga.

\par

Además, también podemos modificar el número de generaciones máximas que queremos que tengan los modelos, la pérdida empleada, el optimizador y las métricas del entrenamiento. En la figura~\ref{fig:codigo_entrenamiento} se muestra la función mencionada con sus parámetros.

\imagen{codigo_entrenamiento}{Función de entrenamiento de los modelos}{1}

Por último, también se pueden modificar las unidades que hay en algunas capas de los modelos, para aumentar o disminuir la complejidad del modelo. Las capas que se pueden modificar tienen un comentario indicándolo en el código.

\par

Tras modificar todo lo deseado podemos ejecutar todas las celdas del fichero y esperar a que se entrenen todos los modelos y ver los resultados en las celdas indicadas para ello.
\apendice{Documentación de usuario}

\section{Introducción}

En este apartado se va a explicar cómo ejecutar el código para entrenar los modelos desarrollados durante el proyecto.

\section{Requisitos de usuarios}

Los requisitos que necesita cumplir el usuario son los siguientes:

\begin{itemize}
    \item 850MB de espacio en el disco duro
    \item Python 3.10.9
    \item Anaconda
\end{itemize}

Se puede encontrar información de cómo instalar Anaconda en su página web~\cite{anaconda}.

\section{Instalación}

Para instalar el código del proyecto, podemos hacerlo desde la siguiente página citada~\cite{github:repo}. Para descargar la carpeta pulsamos el botón \textit{Code} y a continuación pulsamos \textit{Download ZIP}. En la siguiente figura~\ref{fig:github_descarga} podemos ver la ubicación de los botones.

\imagen{github_descarga}{Botones para descargar el repositorio}{1}

\section{Manual del usuario}

Para ejecutar el código que entrena los modelos debemos seguir los siguientes pasos:

\begin{itemize}
    \item Abrimos el fichero \texttt{Predicciones.ipynb}.
    \item Modificamos en la celda de <<Importación y preparación del conjunto de datos>> la ruta de la variable <<path>> para que lea el fichero CSV que contenga los datos que queramos emplear.
    \item Ejecutamos todas las celdas del fichero.
\end{itemize}

A continuación, tras esperar a que se entrenen los modelos, podemos ver en las celdas de rendimiento el error devuelto por cada modelo de forma gráfica y numérica.

Si queremos modificar los modelos, se recomienda leer la documentación técnica de programación~\ref{Documentacion_programacion}, que explica más a detalle el funcionamiento y cómo modificar distintas partes del código.
\setcounter{secnumdepth}{4}

\apendice{Desarrollo del proyecto}

\section{Datos obtenidos de los sensores}

Los datos que se han recogido a través de los sensores se han almacenado en archivos CSV independientes para cada sensor. Se han recogido los siguientes datos cada cinco minutos:

\begin{itemize}
    \item ts (TimeStamp)
    \item fecha
    \item bateria
    \item \texttt{t\_ext} (temperatura exterior (ºC))
    \item \texttt{h\_ext} (\% humedad exterior)
    \item \texttt{t\_C\_cal} (temperatura a poca profundidad calibrada (ºC))
    \item \texttt{h\_C\_cal} (\% humedad a poca profundida calibrada)
    \item \texttt{t\_L\_cal} (temperatura a más profundidad calibrada (ºC))
    \item \texttt{h\_L\_cal} (\% humedad a más profundidad calibrada)
    \item \texttt{h\_C} (\% humedad a poca profundidad)
    \item \texttt{h\_L} (\% humedad a más profundidad)
\end{itemize}

\section{Obtención de datos meteorológicos a través de una API}

Este proyecto comienza con la obtención de los datos meteorológicos de la zona de los sensores para poder validar los datos obtenidos por estos últimos. Para ello, se ha hecho uso de Weatherbit para obtener datos meteorológicos, pero al observar que su límite de datos históricos con la licencia proporcionada para investigación era de un año, se optó por usar OpenWeatherMap, puesto que con la licencia que proporcionaba esta API para estudiantes permitía obtener los datos históricos de más de un año y ha permitido obtener datos históricos correspondientes al dato más antiguo proporcionado por los sensores.

\par

Para obtener los datos, se han realizado llamadas a la API de OpenWeatherMap introduciendo las fechas deseadas y las coordenadas. Se han hecho las llamadas necesarias para obtener datos de la localización de cada sensor y de todo el periodo en el que se han obtenido datos de los sensores.

\par

Tras haber realizado las llamadas y obtener los datos en un formato JSON, se han transformado los datos y se han almacenado en archivos CSV individuales para cada sensor.

\section{Preprocesamiento de datos}

Una vez se han obtenido los datos históricos de cada sensor, se ha realizado un preprocesamiento de los datos.

\par

Para empezar, se ha comprobado la validez de cada dato de los sensores comparándolos con los obtenidos a través de la API. Para comprobar la validez de los datos, se han creado columnas de validez correspondientes a cada uno de los atributos útiles, que serían la temperatura y la humedad exterior y la calibrada, creando un total de 6 columnas nuevas en el conjunto de datos.

\par

Una vez creadas las columnas, se ha empleado una codificación ternaria para la validación de los datos, asignando un valor de 0 a los datos sin validar, un valor de 1 a los datos no válidos y un valor de 3 a los datos válidos.

\par

Para comprobar la validez de los datos, se han comparado los valores recogidos por los sensores con los obtenidos a través de la API. Si el valor absoluto de la diferencia entre los valores supera un margen definido previamente, se considerará un dato no válido. Los márgenes definidos para cada uno de los atributos son los siguientes:

\begin{itemize}
    \item \texttt{t\_ext}: 7ºC
    \item \texttt{h\_ext}: 15\% 
    \item \texttt{t\_C\_cal}: 10ºC
    \item \texttt{h\_C\_cal}: 30\%
    \item \texttt{t\_L\_cal}: 10ºC
    \item \texttt{h\_L\_cal}: 30\%
\end{itemize}

\par

Al realizar la validación de los datos, se ha observado cómo la mayoría de los datos se consideraron no válidos, esto puede suceder debido a que los datos recogidos por los sensores pueden ser distintos a los que se han recogido por las estaciones meteorológicas que emplean las APIs. Esta diferencia puede estar causada por múltiples motivos, como por ejemplo, la composición del suelo donde esa desplegado el sensor, si hay árboles que proyecten sombra sobre el sensor, si hay vegetación alrededor del sensor que suavicen las temperaturas, la humedad puede diferir si se riega el terreno, etc. Por lo que se optó por cambiar el método de validación.

\par

El nuevo método empleado para la validación de los datos fue comparar cada uno de los datos de los sensores con la media móvil de los datos de todos los sensores. La media móvil se compone de los datos de los 15 minutos anteriores y los 15 posteriores en el caso de los atributos exteriores. Para los atributos de temperatura y humedad bajo tierra, se han empleado los datos de los 30 minutos anteriores y los 30 posteriores de todos los sensores.

\imagen{datos_sensor1}{Visualización de la validez de los datos del sensor 1}{0.55}

\par

En el caso del pluviómetro, sí que se han empleado los datos obtenidos a través de la estación meteorológica usada por OpenWeatherMap, puesto que solo disponemos de un pluviómetro desplegado en el viñedo y no podemos realizar comparaciones con otros datos que no sean los históricos.

\par

Puesto que el pluviómetro puede recoger datos erróneos debido a un fuerte viento u otras causas, se ha realizado la validación con un margen de 10 milímetros de lluvia en la variación del pluviómetro.

\imagen{datos_pluv}{Visualización de la validez de los datos del pluviómetro}{0.6}

\par

También se han añadido columnas correspondientes a la variación de la temperatura y la humedad bajo tierra del instante actual respecto al anterior y una columna indicando el mes en el que se ha recogido el dato.

\imagen{datos_columnas}{Dataframe de los datos con todas las columnas}{1}

\section{Visualización previa en MATLAB}

Una vez concluido el preprocesamiento de los datos, se realizaron pruebas previas al desarrollo de las redes neuronales. Estas pruebas consistieron en pruebas de modelos, como una regresión lineal o una predicción de series temporales.

\par

Empleando el modelo de red neuronal proporcionado por MATLAB, se ha realizado una visualización previa de los resultados deseados de la red neuronal a desarrollar. El modelo proporcionado es una red neuronal no linear autorregresiva. Al entrenar este modelo con los datos de uno de los sensores devolvió un MSE de 0.2, podemos usar este valor para establecer un objetivo base.

\section{Preparación de los datos}

Antes de empezar a crear y entrenar modelos, deberemos preparar el conjunto de datos. Para ello, vamos a reducir el conjunto de datos. Puesto que el objetivo es realizar predicciones individuales para cada sensor, solo cogemos los datos de uno de los sensores, en este caso, se han escogido los datos del sensor 3, ya que al compararlo con los demás sensores, se observa que es el que tiene mayor continuidad en los datos.

\par

Tras haber escogido el conjunto de datos a emplear, vamos a reducir más aún la cantidad de datos a emplear. Para empezar, se ha transformado el conjunto de datos para que los datos no sean cada 5 minutos, sino que correspondan a un periodo de una hora, pues queremos realizar predicciones horarias y si lo hiciéramos con un conjunto de datos donde el periodo sea de 5 minutos tendríamos que realizar 12 predicciones y su precisión sería peor. Además, puesto que son datos meteorológicos, al recoger datos cada 5 minutos no se observa apenas diferencia y entraría dentro del rango de error del propio sensor.

\par

Para realizar la transformación del periodo de 5 minutos al periodo horario, se han agrupado los datos en grupos de 12 y se ha almacenado la media de cada uno de los atributos de cada uno de los grupos.

\par

Además, ya que los sensores pueden llegar a funcionar de manera incorrecta y no enviar los datos, pueden llegar a aparecer saltos y discontinuidades entre los datos. Para que eso no afecte en el entrenamiento de los modelos, se ha escogido un periodo en el que no haya ninguna discontinuidad. Ese periodo corresponde a los datos de los últimos 4 meses recogidos por el sensor.

\par

Por último, se han descartado atributos para el entrenamiento de las redes neuronales, puesto que aumentaba la complejidad y aumentaba el MSE de los modelos. Los atributos que se han empleado han sido la temperatura y la humedad bajo tierra, el mes y la hora, un total de 6 atributos.

\section{Normalización}

Para empezar a desarrollar nuestros modelos y comparar su precisión, primero debemos normalizar los valores del conjunto de datos, puesto que facilitará el aprendizaje de nuestros modelos.

\par

Para la normalización de los datos se han probado tanto la puntuación estándar como la puntuación Min-Max. A la hora de ver los resultados, se ha podido comprobar que en este caso particular es mejor la puntuación Min-Max.

\par

Una vez hemos normalizado los datos, dividimos los datos en el conjunto de datos de entrenamiento, el de validación y el de test. Para ello se ha seguido una proporción de 70/20/10 respectivamente. Para dividir el conjunto de datos se ha realizado de manera manual para que las secciones que crearemos a continuación tengan todas las horas seguidas y no haya saltos.

\section{Data windowing}

Para realizar predicciones de series temporales, debemos dividir el conjunto de datos en secciones más pequeñas y entrenar los modelos con esas secciones. Para distintos tipos de predicciones necesitaremos secciones de distintos tamaños.

\par

En este proyecto se han realizado dos tipos de predicciones, predicciones de un solo paso de tiempo y de múltiples pasos de tiempo. Para las predicciones de un solo paso se han creado secciones de 25 datos, donde se introducen como entrada los datos $t_0-t_{23}$ y se quieren predecir los datos $t_1-t_{24}$ de manera individual. En cambio, para las predicciones de múltiples pasos se han creado secciones de 48 datos, donde se introducen los datos $t_0-t_{23}$ y se quieren predecir los datos $t_{24}-t_{47}$.

\par

Para la creación de secciones de datos se ha creado una función específica para ello. Puesto que solo se realizan predicciones de uno de los parámetros del conjunto de datos, introducimos como argumento el parámetro a predecir, que en este caso es \texttt{t\_C\_cal}. Si queremos modificar la variable a predecir, debemos cambiar el argumento a la hora de crear las subsecciones.

\section{Parámetros de entrenamiento}

Para entrenar los modelos se han empleado los mismos parámetros para todos, ya que si no empleamos los mismos no obtendríamos unos resultados que podamos comparar debido a las diferencias en el entrenamiento.

\par

Para ello, se ha creado una función para configurar y entrenar los modelos. La función configura los modelos con los siguientes parámetros:

\begin{itemize}
    \item Nº de generaciones: 200
    \item Paciencia (Nº generaciones sin mejora): 5
    \item Función de pérdida: MSE\footnote{{\label{noteMSE}}Son las siglas del Error Cuadrático Medio en inglés, es una función de pérdida que mide la diferencia entre el valor real y el valor predicho y lo eleva al cuadrado.}
    \item Optimizador: Adam\footnote{Adam\cite{adam} es un optimizador que emplea el algoritmo Adam, que combina dos algoritmos, AdaGrad y RMSProp.}
    \item Métricas: MSE
\end{itemize}

\section{Modelos para predicciones de un solo paso}

Para los modelos para predicciones de un solo paso, es decir, una predicción de un solo dato en el futuro, se ha decidido que las predicciones a realizar sean, teniendo 24 horas de datos como entrada ($t_0-t_{23}$), que realice las predicciones de los instantes siguientes ($t_1-t_{24}$).

\imagen{single_step}{Ejemplos de secciones de datos para las predicciones de un solo paso}{1}

\subsection{Baseline}

Este es el modelo más simple y no necesita de un entrenamiento. Sus predicciones se basan en que no hay cambio, por lo que devuelve el mismo valor que el último dato introducido como entrada.

\par

Podemos observar en la imagen \ref{fig:baseline} que las predicciones corresponden con la hora anterior.

\imagen{baseline}{Resultados del modelo baseline}{0.9}

\subsection{Modelo lineal}

El modelo lineal se compone de una única capa Dense con una unidad, puesto que solo queremos que nos devuelva un atributo como predicción.

\imagen{lineal_modelo}{Capas del modelo lineal}{0.5}

Podemos observar en la siguiente imagen \ref{fig:lineal_single} los resultados.

\imagen{lineal_single}{Resultados del modelo lineal}{0.9}

Se puede ver que las predicciones realizadas por este modelo se alejan bastante del valor real que debería haber predicho.

\subsection{Modelo LSTM}

Por último, para las predicciones de un solo paso también se ha entrenado un modelo LSTM. Este modelo consta de una capa LSTM con 32 unidades y con el parámetro \texttt{return\_sequences=True}, para poder entrenarlo con 24 horas de datos a la vez, y de una capa Dense de una unidad.

\imagen{lstm_modelo}{Capas del modelo LSTM}{0.5}

En la imagen \ref{fig:lstm_single} podemos observar los resultados del modelo, y son bastante acertadas.

\imagen{lstm_single}{Resultados del modelo LSTM}{0.9}

\subsection{Resultados}

Podemos observar en las imágenes~\ref{fig:resultados_single1} y~\ref{fig:resultados_single2} y en la tabla \ref{tabla:res-single} los resultados de los modelos de un paso.

\imagen{resultados_single1}{Resultados de los modelos Baseline y LSTM}{0.7}

\imagen{resultados_single2}{Resultados del modelo lineal}{0.7}

\tablaSmall{Resultados de los modelos de un paso}{ l l l }{res-single}{Modelo & MSE Validación & MSE Test\\}{Baseline & 0.0002 & 0.0003\\
Lineal & 0.0178 & 0.0932\\
LSTM & 0.0002 & 0.0006\\}

Podemos observar que el modelo lineal tiene un MSE muy alto en comparación con los otros dos modelos, teniendo en cuenta que el modelo baseline no realiza ningún cálculo y su predicción se basa en que no hay cambio y usamos este modelo como base para comparar los demás, podemos concluir que las predicciones realizadas por el modelo lineal no son acertadas y se alejan mucho del valor real. Esto se puede deber a múltiples factores, como una falta de entrenamiento o que este tipo de red neuronal no es el mejor para realizar predicciones de series temporales.

\par

En cambio, podemos ver que el modelo LSTM tiene un rendimiento muy bueno a la hora de realizar predicciones de un paso, aunque sigue siendo peor que el baseline, pero el error es tan pequeño que entra dentro del margen de error, pues el error es casi imperceptible.

\section{Modelos para predicciones de múltiples pasos}

Las predicciones de múltiples pasos son predicciones que pueden realizarse varios pasos de tiempo en el futuro. En este caso, se han realizado predicciones de las 24 horas posteriores ($t_{24}-t_{47}$) teniendo como datos de entrada las 24 horas anteriores ($t_0-t_{23}$).

\par

Para estas predicciones se han entrenado los modelos 10 veces y se ha almacenado la media de las ejecuciones para obtener unos resultados más fiables y menos afectadas por la aleatoriedad.

\imagen{multi_step}{Ejemplos de secciones para las predicciones de múltiples pasos}{1}

\subsection{Modelo lineal}

Para las predicciones de múltiples pasos vamos a emplear el modelo lineal como modelo base para comparar el rendimiento de los demás. El modelo lineal realiza todas las predicciones necesarias a partir del último dato de entrada, es decir, la última hora que introducimos.

\par

Este modelo está compuesto de una capa Lambda que solo lee el último dato de entrada, una capa Dense y una capa Reshape para modificar la salida de la capa anterior.

\imagen{lineal_multi_modelo}{Capas del modelo lineal}{0.5}

\par

Podemos ver las predicciones realizadas por este modelo en la imagen \ref{fig:lineal_multi}.

\imagen{lineal_multi}{Predicciones de múltiples pasos del modelo lineal}{1}

Se puede observar cómo la predicción de la primera hora es bastante cercana a su valor real, sin embargo, las siguientes horas se empiezan a alejar del valor que debería predecir. En cambio, se puede ver cómo las predicciones tienen una curva que se asemeja a la variación de la temperatura a lo largo del día, aunque las predicciones sean erróneas.

\subsection{Modelo denso}

El modelo denso es igual que el modelo lineal, solo que añade una capa Dense extra después de la capa Lambda para añadirle más potencia, sin embargo, sigue basando todas las predicciones en el último paso de tiempo de entrada. Esta capa Dense tiene 512 unidades y tiene la función de activación \textit{relu}.

\imagen{dense_modelo}{Capas del modelo denso}{0.5}

Las predicciones realizadas por este modelo se pueden visualizar en la siguiente imagen\ref{fig:dense_multi}.

\imagen{dense_multi}{Predicciones de múltiples pasos del modelo denso}{1}

Todas las predicciones se acercan bastante a su valor real y la curva de la variación de la temperatura se asemeja bastante. Podemos notar comparando estos resultados con los anteriores que la capa adicional sí que ha aumentado la potencia de la red neuronal a pesar de solo basar las predicciones en la última hora.

\subsection{Modelo convolucional}

El modelo convolucional realiza las predicciones en base a un historial de las últimas horas, lo que aumenta su precisión respecto al modelo anterior. En este caso, el modelo convolucional tiene un historial de todas las horas introducidas como entrada.

\par

Este modelo añade una capa adicional al modelo lineal que hemos mencionado antes, sin embargo, esta capa es de tipo \textit{Conv1D} y tiene 256 unidades y tiene un tamaño de kernel de 24, equivalente al tamaño del historial que queremos. Además, modifica la capa Lambda para que sea capaz de leer todas las horas y no solo la última.

\imagen{cnn_modelo}{Capas del modelo convolucional}{0.5}

Podemos ver que las predicciones devueltas por este modelo solo presenta una pequeña mejora frente al modelo anterior y es capaz de captar mejor la curva de la variación de la temperatura a lo largo del día, por lo que también es capaz de realizar predicciones más acertadas.

\imagen{cnn_multi}{Predicciones de múltiples pasos del modelo convolucional}{1}

\subsection{Modelo LSTM}

El modelo LSTM es un modelo recurrente que entrena un modelo de manera recurrente tantas veces como datos de entrada introduzcamos y luego realiza las predicciones con el último modelo.

\par

Este modelo está compuesto por una capa Dense con 32 unidades y con el parámetro \texttt{return\_sequences=False}, una capa Dense y una capa Reshape.

\imagen{lstm_multi_modelo}{Capas del modelo LSTM}{0.5}

Podemos observar que los resultados de este modelo son bastante similares al modelo convolucional, si embargo, se puede apreciar que las primeras predicciones devueltas suelen tener una mayor variación respecto a su valor real, pero las siguientes se acercan más al valor que debería predecir.

\imagen{lstm_multi}{Predicciones de múltiples pasos del modelo LSTM}{0.9}

Se podría concluir que este modelo no es capaz de predecir correctamente las primeras horas a la hora de realizar predicciones de múltiples pasos, sin embargo, su precisión mejora a medida que aumentamos los pasos de tiempo.

\subsection{Modelo LSTM autorregresivo}

Para finalizar, el último modelo entrenado y comparado es un modelo LSTM autorregresivo, que entrena el modelo de la misma manera que el modelo anterior, sin embargo, cada vez que realiza una predicción, vuelve a entrenar el modelo.

\par

Para este modelo se ha creado una clase FeedBack, que representa el modelo, y emplea una capa LSTMCell dentro de una capa RNN y, a continuación de esta, una capa Dense.

\par

A continuación, se ha creado un método de calentamiento para inicializar el estado interno en base a la entrada, que sería equivalente al modelo LSTM para predicciones de un solo paso. Tras la creación de este método, se ha modificado el método \textit{call} del modelo para que inicialice el estado empleando el método creado, haga la primera predicción con el calentamiento y realice el resto de las predicciones.

\par

Podemos ver que las predicciones realizadas por este modelo suele realizar las primeras horas de manera más o menos acertada. Sin embargo, tras unas cuantas predicciones devuelve unos valores anómalos que se alejan bastante de la curva de la variación de la temperatura y del valor real a predecir.

\imagen{lstm_ar_multi}{Predicciones de múltiples pasos del modelo LSTM autorregresivo}{0.9}

\subsection{Resultados}

Podemos observar en la siguiente tabla y en la imagen los valores correspondientes al MSE de cada modelo.

\tablaSmall{Resultados de los modelos de múltiples pasos}{ l l l }{res-multi}{Modelo & MSE Validación & MSE Test\\}{Lineal & 0.0175 & 0.0602\\
Denso & 0.0064 & 0.0324\\
Convolucional & 0.0039 & 0.0224\\
LSTM & 0.0052 & 0.0425\\
LSTM autorregresivo & 0.0056 & 0.0419\\}

\imagen{resultados_multi}{Resultados de los modelos de múltiples pasos}{0.6}

Se puede apreciar cómo los modelos más simples y el más complejo tienen un peor rendimiento respecto a los modelos que se encuentran entre ellos. Ya que los modelos lineal y denso solo realizan las predicciones en base al último paso de tiempo que se introduce como entrada, cabe esperar que sus predicciones tampoco sean tan exactas como las de los demás modelos.

\par

También podemos observar que el modelo LSTM autorregresivo tiene una precisión parecida comparada con su modelo más simple para realizar predicciones de múltiples pasos. Esto se puede deber a que un modelo demasiado complejo disponga de un número de pesos demasiado elevado para el entrenamiento y haga predicciones no tan acertadas y no sigan una tendencia regular.

\par

Cabe mencionar que todos estos modelos no son capaces de realizar una predicción con precisión cuando las temperaturas del día a predecir varían mucho respecto a las del día anterior, puesto que como solo introducimos 24 horas, las predicciones realizadas por los modelos seguirán la tendencia del día anterior y no se saldrán casi de los valores máximos o mínimos del día anterior.

\bibliographystyle{plain}
\bibliography{bibliografiaAnexos}

\end{document}
